% !PRINTED 105
\weberpage{105}
%<<CONT
and obtain, if $\varepsilon_1$, $\varepsilon_2$ have the value 0 or 1, the resolvents

\weq{6}{\begin{aligned}
\Omega_{\varepsilon_1,\varepsilon_2} &= \sum^n (-1)^{\varepsilon_1\sigma_1+\varepsilon_2\sigma_2}\, r^n\\
&= \eta + (-1)^{\varepsilon_1}\eta' + (-1)^{\varepsilon_2}\eta'' + (-1)^{\varepsilon_1+\varepsilon_2}\eta'''.
\end{aligned}}

If now $m = q_1\, q_2 \ldots q_\mu$ is not divisible by 8, then, if $r_1, r_2, \ldots r_\mu$
denote primitive roots of the degrees $q_1, q_2, \ldots q_\mu$, we can
set
\[ r = r_1\, r_2 \ldots r_\mu \]
and further
\[ \begin{aligned}
n &\equiv n_1\ (\mathrm{mod}\ q_1)\\
&\equiv n_2\ (\mathrm{mod}\ q_2)\\
&\cdot\ \cdot\ \cdot\ \cdot\ \cdot\ \cdot\ \cdot\\
&\equiv n_\mu\ (\mathrm{mod}\ q_\mu),
\end{aligned} \]
from which it results:
\weq{7}{\Omega_{\varepsilon_1,\varepsilon_2} = \prod_{1,\mu}^{h} \sum^{n_h} (-1)^{\nu_h(\varepsilon_1 e_{1,h}+\varepsilon_2 e_{2,h})}\, r_h^{n_h}.}

The individual factors of this product are the \gesperrt{Gauss sums}
determined in \S.\,171 (Bd.\ I), which are there expressed by
square roots. Only when $q_1 = 4$ and
$\varepsilon_1\, e_{1,1} + \varepsilon_2\, e_{1,2} = 0$ is $\Omega_{\varepsilon_1,\varepsilon_2} = 0$. % [sic: Weber prints e_{1,2}; specializing the preceding general factor to h=1 requires e_{2,1}.]

But if $m$ is divisible by 8 and
\[ m = q_2\, q_3 \ldots q_\mu, \quad q_2 = 8, \]
then, setting $r = r_2\, r_3 \ldots r_\mu$, one obtains:
\[ \begin{aligned}
\Omega_{\varepsilon_1,\varepsilon_2} &= \sum^{n_2} (-1)^{\varepsilon_1(e_{1,1}\nu_1+e_{1,2}\nu_2)}(-1)^{\varepsilon_2(e_{2,1}\nu_1+e_{2,2}\nu_2)}\, r_2^{n_2}\\
&\times \prod_{3,\mu}^{h} \sum^{n_h} (-1)^{\nu_h(\varepsilon_1 e_{1,h}+\varepsilon_2 e_{2,h})}\, r_h^{n_h}.
\end{aligned} \]

In the five groups set up above, which belong to $m = 24$,
the periods can be computed from this formula, or also
easily directly, by setting
\[ r = e^{\frac{\pi i}{12}} = e^{\frac{\pi i}{3}}\, e^{-\frac{\pi i}{4}} \]
and one then obtains for the five quantities:
\[ \mathfrak{A}_1)\quad \eta = \frac{(1+i)\sqrt{3}}{\sqrt{2}} \]
\[ \mathfrak{A}_2)\quad \eta = \frac{\sqrt{3}-i}{\sqrt{2}} \]
%CONT>>
